% Copyright (c) 2026 Eugene Joseph M. Ragasa.
% SPDX-License-Identifier: Apache-2.0
% Author: Eugene Joseph M. Ragasa
% Formal proof completed: 2026-10-19
\documentclass[11pt]{article}
\usepackage[T1]{fontenc}
\usepackage{amsmath,amssymb,booktabs,tabularx}
\usepackage[margin=1in]{geometry}
\title{Proof Dossier: Hermitian-Field Fourier Pairing}
\author{Eugene Joseph M. Ragasa}
\date{Formal proof completed: 2026-10-19}
\begin{document}\maketitle
\noindent\begin{tabularx}{\textwidth}{@{}lX@{}}\toprule
Statement ID & \texttt{BG-FOURIER-THM-022} \\
Status & proved; no admissions \\
Lean declaration & \texttt{ProjectKoios.BandGauge.Fourier.}\newline\texttt{finiteFourierCoefficient\_neg\_eq\_conjTranspose} \\
Owning module & \texttt{ProjectKoios/BandGauge/Fourier/HermitianPairing.lean} \\
Axioms & \texttt{propext}, \texttt{Classical.choice}, \texttt{Quot.sound} \\
\bottomrule\end{tabularx}
\section{Statement and argument}
For a pointwise Hermitian finite matrix field, the coefficient at the negative character equals the conjugate transpose of the coefficient at the original character. Expand entries, conjugate the finite sum, and use Hermiticity pointwise.
\section{Boundary}
This does not bind characters to physical lattice vectors or a publication mesh. Lean checked the exact statement; the dossier and proof were AI-assisted.
\end{document}
